% Source: https://artofproblemsolving.com/community/c3h2295883_prime_factorization_problem_with_unknowns
% Stable identity: aops-topic-2295883-q01
The number $7+30N$ leaves remainder $7$ when divided by $30$, so it is divisible by none of $2$, $3$, $5$. For $N\le5$ we have $7+30N\le157<169=13^2$, so if such a number were composite it would have a prime factor less than $13$ other than $2,3,5$, that is, $7$ or $11$.
\begin{itemize}
\item Modulo $7$: $7+30N\equiv2N$, which is divisible by $7$ only when $7\mid N$.
\item Modulo $11$: $7+30N\equiv7+8N$, which is divisible by $11$ exactly when $8N\equiv4$, i.e. $2N\equiv1$, i.e. $N\equiv6\pmod{11}$.
\end{itemize}
So $7+30N$ is prime for $N=1,2,3,4,5$ (the numbers $37,67,97,127,157$), while for $N=6$ it is $187=11\cdot17$. The answer is $\boxed{N=6}$.
